\section{一般四窓の切断と七次までの共通Tor検査}
\label{sec:general-cut-trivial-extension}
三因子対称性を使わずにも、戻り辺の切断は有限Frobenius候補全体に適用できる。
ここでは有限次数での必要条件を整理し、全Koszul性の有限判定とは区別する。

\begin{proposition}[任意の一辺切断によるtrivial extension]
\label{prop:general-cut-trivial-extension}
標準Frobenius四窓を折り畳んだ代数$E(w)$で、次数一の四辺のうち一辺を選ぶ。
その辺を通らない経路の部分代数を$A$とすると、$A$は32次元の四層二次Koszul代数で、
大域次元は3以下である。また
\[
 E(w)\simeq A\ltimes DA
\]
が成り立つ。$A$の長さ$r$成分の双対には次数$4-r$を入れる。
四つの切断のいずれでもこの表示を取れる。
これは一般対象に可換性・domain・noetherian性を加える主張ではない。
\end{proposition}
\begin{proof}
切断後の頂点を$0,1,2,3$の順に置く。
$A$の次数別次元は$4,12,12,4$であり、全次元は32である。
端点と長さを固定した経路は切断辺を通る回数が一定なので、
二次関係を課しても「切断辺を通らない部分」と「一回通る部分」の直和は保たれる。
前者は四層の二次表示を持つ部分代数であり、補題\ref{lem:four-layer-koszul}を適用できる。
後者を$M$と書くと$E=A\oplus M$で、$M$は双加群イデアルである。
切断辺を二回通る経路の長さは5以上であり、$E$の五次以上は零なので$M^2=0$である。

標準traceの巡回恒等式から、四頂点のtraceの和$\lambda$は対称なFrobenius形式を与える。
$\lambda(AA)=0$、$\lambda(MM)=0$なので、完全対合は$A$と$M$を完全に対合させる。
写像$m\mapsto(a\mapsto\lambda(am))$は$M\simeq DA$を与え、
traceの巡回性から左右の$A$作用を保つ。次数は$\lambda$の次数4から決まる。
従って平方零拡大の積は通常のtrivial extensionの積そのものである。
\end{proof}

切断後の位置$i,j\in\{0,1,2,3\}$を固定し、$d=4+j-i$と置く。
左右のKoszul分解を使って命題\ref{prop:one-return-bar-tor}の一回部分を計算する。
$h$番目の項$K_h^{i,j}$では、左右のKoszul次数を$p,q$として
\[
 p+q=h,\qquad 0\leq p\leq3-i,\quad0\leq q\leq j,
 \qquad r=i-j+h
\]
となり、中央係数は長さ$r$の$A$成分の双対である。
$h_r=(1,4,6,4)_r$とし、範囲外では$h_r=0$と定めれば
\[
 \dim K_h^{i,j}
 =h_{i-j+h}\!!\sum_{\substack{p+q=h\\p\leq3-i,\ q\leq j}}
       \binom{p+3}{3}\binom{q+3}{3}.
\]
ここでKoszul空間の次元$1,4,10,20$は、$A$のHilbert行列の逆行列から従う。
異なる窓の空間を同一視することは仮定していない。

\begin{proposition}[七次切断の必要十分な共通検査]
\label{prop:cut-tor-through-seven}
四つの切断それぞれで、全16組$(i,j)$の上の複体が最上次数$d-1$より下で正確であることは、
$E(w)$の内部次数$1\leq d\leq7$で$\operatorname{Tor}_{n,d}^{E(w)}=0$ ($n\ne d$)となることと同値である。
正確な場合の最上ホモロジー次元は$\binom{d+3}{3}$である。
特に$d=7$の一回部分は、零でない項の次元が
\[
 (\dim K_3,\dim K_4,\dim K_5,\dim K_6)=(120,1040,2400,1600)
\]
となる小複体で検査できる。
この有限検査だけから全Koszul性やAS正則性は結論しない。
\end{proposition}
\begin{proof}
一回部分の内部次数は$4+j-i$、最上Koszul次数は$3+j-i=d-1$である。
そのホモロジーは$E$の$\operatorname{Tor}_{h+1,d}$の対応する直和因子となる。
長さ7以下の有向四周期経路では、少なくとも一つの辺をちょうど一回通る。
その辺で切れば、端点と内部次数を固定した被約バー複体全体が戻り辺一回部分になる。
従って四切断の検査は内部次数7以下の全Tor成分を覆い、述べた同値を得る。
最上次数の符号を正に取った上の次元公式のEuler標数は
$\sum_h(-1)^{d-1-h}\dim K_h^{i,j}=\binom{d+3}{3}$である。
これは$0\leq i,j\leq3$の有限和で直接計算でき、$d=7$の表示も同じ公式による。

長さ8では全ての辺を二回通る経路があり、一回部分では覆えない。
また、より高次の非線型Torの消滅はここでは証明していないので、全Koszul性の十分条件とはしない。
\end{proof}

\path{computations/verify_cut_tor_dimensions.py}は全16組の項次元とEuler標数、
長さ7までの全始点の経路に一回辺があることを整数で検算する。
これは行列の階数や実際の複体の正確性を検算するものではない。

\paragraph{二つの有理点での全微分。}
\path{computations/general_cut_tor.py}では二次関係を挿入した行列の核から
各$Q_p$と中央の双対成分の基底を取る。左作用は左因子と中央因子の末尾の文字同士を縮約する。
右作用は両者の先頭を縮約し、符号$(-1)^p$を付ける。
窓の位置をこの二方向で別々に追うので、中央に可換性を仮定しない。
有理数上で四切断各16組の全ての連続微分の合成と階数を計算した。
符号付きOre例$H=0$, $G=(1,0,0)$では全始点、内部次数7以下で非対角成分は零であり、
特に内部次数7の三階数は小さい$h$から$120,920,1480$となる。
一方、節\ref{sec:trivial-extension-bitor}のBoolean境界$t=0$で非対角成分が出る位置は次の通りである。
表は$\dim\operatorname{Tor}^{E}_{d-1,d}$であり、四つの列は元の始点に対応する。
\[
\begin{array}{c|rrrr}
d&0&1&2&3\\ \hline
5&7&0&0&0\\
6&21&0&0&21\\
7&41&0&41&63
\end{array}
\]
この範囲で他の非対角成分は零、各始点の対角成分の次元は表の値に$\binom{d+3}{3}$を足した数になる。
同じ始点・次数を覆う切断は全て同じホモロジーを与えた。計算入力、再現手順と階数は
\path{computations/verify_general_cut_tor.py}および\path{research/general_cut_tor_certificate.json}に置く。
